\section{Yadroning bog‘lanish energiyasi}



\begin{enumerate}

\item Quyida keltirilgan zarralarning qaysilari orasida yadro kuchlari mavjud?
1) proton; 2) neytron; 3) elektron; 4) pozitron; 5) neytrino; 6) $\alpha$-zarra.
\begin{enumerate}[label=\Alph*)]
\item 1, 2.
\item 3, 4.
\item 1, 3.
\item 2, 5.
\end{enumerate}

\item Tinch turgan elektron va proton bir-biri bilan qanday kuchlar vositasida ta’sirlashadi?
1) Kulon kuchi; 2) Amper kuchi; 3) Lorens kuchi; 4) yadro kuchlari; 5) gravitatsiya kuchi.
\begin{enumerate}[label=\Alph*)]
\item 1 va 5.
\item 1 va 3.
\item 1, 2 va 5.
\item 1, 3 va 5.
\end{enumerate}

\item Atom yadrosi mustahkam bo‘lishi uchun ...
\begin{itemize}
    \item[I.] Protonlari soni kam bo‘lishi kerak.
    \item[II.] Radioaktiv bo‘lishi kerak.
    \item[III.] Bog‘lanish energiyasi katta bo‘lishi kerak.
\end{itemize}
\begin{enumerate}[label=\Alph*)]
\item faqat I.
\item faqat II.
\item faqat III.
\item I va III.
\end{enumerate}

\item Quyidagi fikrlardan qaysilari to‘g‘ri?
\begin{itemize}
    \item[I.] Elektron atomning markazida joylashgan.
    \item[II.] Neytral atomda elektronlar soni protonlar soniga teng.
    \item[III.] Yadroni tashkil qiluvchi zarralar orasida juda qisqa masofada yadroviy kuchlar bor.
\end{itemize}
\begin{enumerate}[label=\Alph*)]
\item faqat I.
\item I va II.
\item I va III.
\item II va III.
\end{enumerate}

\setcounter{enumi}{4}
\item 3 ta proton va 4 ta neytrondan iborat yadroning solishtirma bog‘lanish energiyasi 5,6 MeV/nukl. Shu yadroni erkin proton va neytronlarga ajratib yuborish uchun kamida qancha energiya zarur bo‘ladi (MeV)?
\begin{enumerate}[label=\Alph*)]
\item 0,8.
\item 5,6.
\item 11,2.
\item 39,2.
\end{enumerate}

\item Litiy $^7_3\text{Li}$ yadrosidagi nuklonlarning solishtirma bog‘lanish energiyasi 5,4 MeV/nuklon. Yadroning massa deffektini (kg) toping. $1\text{ MeV} = 1,6\cdot 10^{-13}$ J.
\begin{enumerate}[label=\Alph*)]
\item $6,05\cdot 10^{-14}$.
\item $0,96\cdot 10^{-29}$.
\item $6,05\cdot 10^{-12}$.
\item $6,72\cdot 10^{-29}$.
\end{enumerate}

\item Tinchlikdagi massalari $m_0$ dan bo‘lgan ikki zarra birlashganda E energiya ajraladi. Hosil bo‘lgan zarraning tinchlikdagi massasi qanday?
\begin{enumerate}[label=\Alph*)]
\item $2m_0 + E/c^2$.
\item $2m_0c^2 + E$.
\item $2m_0 - E/c^2$.
\item $2m_0 - Ec^2$.
\end{enumerate}

\item Aluminiy $^{27}_{13}\text{Al}$ yadrosining bog‘lanish energiyasi ifodasini ko‘rsating ($m_p$ -- proton massasi, $m_n$ -- neytron massasi, $M_{ya}$ -- yadro massasi).
\begin{enumerate}[label=\Alph*)]
\item $M_{ya}c^2$.
\item $(13m_p - 14m_n - M_{ya})c^2$.
\item $(27m_p + 13m_n - M_{ya})c^2$.
\item $(13m_p + 14m_n - M_{ya})c^2$.
\end{enumerate}

\item Protonlar soni Z, neytronlar soni N bo‘lgan yadroning bog‘lanish energiyasining matematik ifodasi qaysi javobda to‘g‘ri berilgan? $m_{ya}$ -- yadro massasi, $m_p$ -- proton massasi, $m_n$ -- neytron massasi.
\begin{enumerate}[label=\Alph*)]
\item $m_{ya}c^2$.
\item $(m_{ya} + Zm_p + Nm_n)c^2$.
\item $(m_{ya} - Zm_p - Nm_n)c^2$.
\item $(Zm_p + Nm_n - m_{ya})c^2$.
\end{enumerate}

\setcounter{enumi}{9}
\item Geliy atomi yadrosining bog‘lanish energiyasi $4,5\cdot 10^{-12}$ J bo‘lsa, uning massa defekti qanday (kg)?
\begin{enumerate}[label=\Alph*)]
\item $1,5\cdot 10^{-20}$.
\item $2\cdot 10^{-29}$.
\item $5\cdot 10^{-29}$.
\item $2,25\cdot 10^{-29}$.
\end{enumerate}

\item 44 ta proton va 56 ta neytrondan iborat yadroning bog‘lanish energiyasi 800 MeV bo‘lsa, yadroning solishtirma bog‘lanish energiyasi nimaga teng bo‘ladi (MeV)?
\begin{enumerate}[label=\Alph*)]
\item 700.
\item 67.
\item 18.
\item 8.
\end{enumerate}

\item $^7_3\text{Li}$ izotopi massa defekti $6,89\cdot 10^{-26}$ g. Xuddi shunday izotopning 7 g massasini alohida proton va neytronlarga ajratib yuborish uchun qancha energiya (MJ) kerak bo‘ladi?
\begin{enumerate}[label=\Alph*)]
\item $1,86\cdot 10^6$.
\item $3,72\cdot 10^{12}$.
\item $3,72\cdot 10^6$.
\item $1,86\cdot 10^{12}$.
\end{enumerate}

\end{enumerate}





\begin{table}[!ht]
\centering
\begin{tabular}{|c|c|c|c|c|c|c|c|c|c|c|}
\hline
1a & 2a & 3d & 4d & 5d & 6d & 7c & 8d & 9d & 10c & 11d \\ \hline
12b & & & & & & & & & & \\ \hline
\end{tabular}
\end{table}